% GENERATED FILE. Edit canonical lessons, then run npm run generate:books.
% canonical-source-hash: fnv1a64:c68e88ea9cdee622
% canonical-lessons: TE-C263-verusanaga, TE-C263-ruci, TE-C263-pyantu, TE-C263-nikkaru, TE-C263-baninu

\chapter{Peanuts, Taste, Trousers, Shorts, Vest}
\label{ch:te-peanuts-taste-trousers-shorts-vest}
\hlchaptermodality{263}

\begin{chapteropening}
\textbf{By the end of this chapter:} \emph{I can say peanuts, taste, trousers, shorts and a vest.}

Complete the last lesson of chapter 263: I can say peanuts, taste, trousers, shorts and a vest.
\end{chapteropening}

\section[\texorpdfstring{\te{వేరుశనగ} (vēruśanaga)}{వేరుశనగ (vēruśanaga)}]{\te{వేరుశనగ} (vēruśanaga) — peanuts}

\label{lesson:TE-C263-verusanaga}



\begin{quote}
Before the new one: say the Telugu for breakfast, then the Telugu for green gram.
\end{quote}

\subsection*{You'll want to know: \te{వేరుశనగ}}
\textbf{\te{వేరుశనగ}} — \emph{vēruśanaga} — \textquotedblleft{}peanuts\textquotedblright{}.

\begin{grammarlens}[title={What you've built}]
One more word for this chapter.
\end{grammarlens}

\subsection*{Your turn}
\begin{itemize}
\raggedright
  \item \emph{Say it:} \emph{vēruśanaga}
  \item \emph{Say it:} \emph{vēruśanaga}, once more
  \item \emph{Say it:} say \emph{vēruśanaga}
  \item \emph{Recall:} say the Telugu for breakfast, then the Telugu for green gram, then say \emph{vēruśanaga} again
\end{itemize}

\begin{tcolorbox}[breakable,colback=teal!4,colframe=teal!35!black,title={Before you move on}]
What does \te{వేరుశనగ} mean? (\textquotedblleft{}Peanuts\textquotedblright{}.) Say it once more.
\end{tcolorbox}
\section[\texorpdfstring{\te{రుచి} (ruci)}{రుచి (ruci)}]{\te{రుచి} (ruci) — taste, flavour}

\label{lesson:TE-C263-ruci}



\begin{quote}
Before the new one: say the Telugu for green gram, then the Telugu for peanuts.
\end{quote}

\subsection*{You'll want to know: \te{రుచి}}
\textbf{\te{రుచి}} — \emph{ruci} — \textquotedblleft{}taste, flavour\textquotedblright{}.

\begin{grammarlens}[title={What you've built}]
One more word for this chapter.
\end{grammarlens}

\subsection*{Your turn}
\begin{itemize}
\raggedright
  \item \emph{Say it:} \emph{ruci}
  \item \emph{Say it:} \emph{ruci}, once more
  \item \emph{Say it:} say \emph{ruci}
  \item \emph{Recall:} say the Telugu for green gram, then the Telugu for peanuts, then say \emph{ruci} again
\end{itemize}

\begin{tcolorbox}[breakable,colback=teal!4,colframe=teal!35!black,title={Before you move on}]
What does \te{రుచి} mean? (\textquotedblleft{}Taste, flavour\textquotedblright{}.) Say it once more.
\end{tcolorbox}
\section[\texorpdfstring{\te{ప్యాంటు} (pyāṇṭu)}{ప్యాంటు (pyāṇṭu)}]{\te{ప్యాంటు} (pyāṇṭu) — trousers}

\label{lesson:TE-C263-pyantu}



\begin{quote}
Before the new one: say the Telugu for peanuts, then the Telugu for taste.
\end{quote}

\subsection*{You'll want to know: \te{ప్యాంటు}}
\textbf{\te{ప్యాంటు}} — \emph{pyāṇṭu} — \textquotedblleft{}trousers\textquotedblright{}.

\begin{grammarlens}[title={What you've built}]
One more word for this chapter.
\end{grammarlens}

\subsection*{Your turn}
\begin{itemize}
\raggedright
  \item \emph{Say it:} \emph{pyāṇṭu}
  \item \emph{Say it:} \emph{pyāṇṭu}, once more
  \item \emph{Say it:} say \emph{pyāṇṭu}
  \item \emph{Recall:} say the Telugu for peanuts, then the Telugu for taste, then say \emph{pyāṇṭu} again
\end{itemize}

\begin{tcolorbox}[breakable,colback=teal!4,colframe=teal!35!black,title={Before you move on}]
What does \te{ప్యాంటు} mean? (\textquotedblleft{}Trousers\textquotedblright{}.) Say it once more.
\end{tcolorbox}
\section[\texorpdfstring{\te{నిక్కరు} (nikkaru)}{నిక్కరు (nikkaru)}]{\te{నిక్కరు} (nikkaru) — shorts}

\label{lesson:TE-C263-nikkaru}



\begin{quote}
Before the new one: say the Telugu for taste, then the Telugu for trousers.
\end{quote}

\subsection*{You'll want to know: \te{నిక్కరు}}
\textbf{\te{నిక్కరు}} — \emph{nikkaru} — \textquotedblleft{}shorts\textquotedblright{}.

\begin{grammarlens}[title={What you've built}]
One more word for this chapter.
\end{grammarlens}

\subsection*{Your turn}
\begin{itemize}
\raggedright
  \item \emph{Say it:} \emph{nikkaru}
  \item \emph{Say it:} \emph{nikkaru}, once more
  \item \emph{Say it:} say \emph{nikkaru}
  \item \emph{Recall:} say the Telugu for taste, then the Telugu for trousers, then say \emph{nikkaru} again
\end{itemize}

\begin{tcolorbox}[breakable,colback=teal!4,colframe=teal!35!black,title={Before you move on}]
What does \te{నిక్కరు} mean? (\textquotedblleft{}Shorts\textquotedblright{}.) Say it once more.
\end{tcolorbox}
\section[\texorpdfstring{\te{బనీను} (banīnu)}{బనీను (banīnu)}]{\te{బనీను} (banīnu) — a vest, an undershirt}

\label{lesson:TE-C263-baninu}



\begin{quote}
Before the new one: say the Telugu for trousers, then the Telugu for shorts.
\end{quote}

\subsection*{You'll want to know: \te{బనీను}}
\textbf{\te{బనీను}} — \emph{banīnu} — \textquotedblleft{}a vest, an undershirt\textquotedblright{}.

\begin{grammarlens}[title={What you've built}]
One more word for this chapter.
\end{grammarlens}

\subsection*{Your turn}
\begin{itemize}
\raggedright
  \item \emph{Say it:} \emph{banīnu}
  \item \emph{Say it:} \emph{banīnu}, once more
  \item \emph{Say it:} say \emph{banīnu}
  \item \emph{Recall:} say the Telugu for trousers, then the Telugu for shorts, then say \emph{banīnu} again
\end{itemize}

\begin{tcolorbox}[breakable,colback=teal!4,colframe=teal!35!black,title={Before you move on}]
What does \te{బనీను} mean? (\textquotedblleft{}A vest, an undershirt\textquotedblright{}.) Say it once more.
\end{tcolorbox}
